\section{The Compactness-Partisan Correlation Within the Ensemble}
\label{sec:correlation}

\subsection{The Rodden Effect in Ensemble Plans}

\citet{rodden2019why} documented that geographic concentration of Democratic
voters in cities creates a structural Republican advantage under compact
districting.  This effect is visible within \recom{} ensembles: sorting
ensemble plans by compactness reveals a correlation between higher compactness
and fewer Democratic seats in sorted states.

Let $\rho(\pp, S_D)$ denote the Pearson correlation between mean Polsby-Popper
score and Democratic seat count across ensemble plans.  For the states studied:

\begin{center}
\begin{tabular}{lcc}
\toprule
State & $\rho(\pp, S_D)$ & Direction \\
\midrule
NC & $-0.18$ & Higher PP $\to$ slightly fewer D seats \\
WI & $-0.09$ & Weak (competitive) \\
GA & $-0.31$ & Higher PP $\to$ fewer D seats \\
PA & $-0.11$ & Weak (competitive) \\
\bottomrule
\end{tabular}
\end{center}

The correlation is negative in all states: more compact plans tend to produce
fewer Democratic seats.  The effect is stronger in sorted states (Georgia)
and weaker in competitive states (Wisconsin, Pennsylvania).

\subsection{Implications for AR's Position}

The \ar{} plan is at the $61$st--$72$nd percentile of compactness.  The
negative correlation between compactness and Democratic seats implies that
a plan at this percentile will tend to have slightly fewer Democratic seats
than the ensemble median.

The correct regression formula for the expected seat-count change is:
\begin{equation}
  \Delta S_D \approx \rho \cdot \frac{\sigma_{S_D}}{\sigma_\pp} \cdot \Delta\bar\pp,
  \label{eq:delta-seats}
\end{equation}
where $\Delta\bar\pp = \bar\pp_{\ar} - \mu_{\pp,\rm ens}$ is the absolute
deviation of the \ar{} plan's PP score from the ensemble mean.

For North Carolina, $\rho = -0.18$ and
$\Delta\bar\pp = 0.337 - 0.318 = 0.019$:

\begin{equation}
  \Delta S_D \approx (-0.18) \cdot \frac{1.1}{0.031} \cdot 0.019 \approx -0.12 \text{ seats},
\end{equation}

which rounds to zero.  The NC correlation is too weak to produce a
measurable seat change at this percentile level.

For Georgia, $\rho = -0.31$,
$\Delta\bar\pp = 0.395 - 0.312 = 0.083$:

\begin{equation}
  \Delta S_D \approx (-0.31) \cdot \frac{1.2}{0.030} \cdot 0.083 \approx -1.03 \text{ seats},
\end{equation}

which is essentially exactly the observed $-1$ seat deviation (5D vs.\ median
6D).  The corrected formula thus quantitatively \emph{explains} the full Georgia
deviation via the compactness-partisan correlation alone, with no residual to
attribute to the partition structure of $k = 14 = 2 \times 7$.  This
strengthens the paper's central argument: the Georgia outcome is a geometric
consequence of the Rodden sorting mechanism, not an artefact of the algorithm.

\subsection{Is the Correlation Causal?}

The negative $\rho(\pp, S_D)$ is observational, not causal.  Compact plans
do not cause fewer Democratic seats; rather, the same geographic property
(urban concentration of Democrats) that makes compact plans produce Republican
advantages also determines the correlation within the ensemble.

The policy implication: one cannot avoid the Rodden effect by choosing a
different algorithm.  Any compact algorithm applied to a sorted state will
produce a similar correlation.  The correlation is a property of geography,
not of the \ar{} algorithm.
